% GENERATED FILE. Edit canonical lessons, then run npm run generate:books.
% canonical-source-hash: fnv1a64:c770b967873a8446
% canonical-lessons: PA-C303-nahir, PA-C303-khuh, PA-C303-sinjna, PA-C303-kisan, PA-C303-halat

\chapter{Water for the Fields}
\label{ch:pa-a2-water-for-the-fields}
\hlchaptermodality{303}

\begin{chapteropening}
\textbf{By the end of this chapter:} \emph{I can say how the fields get their water.}

Complete the last lesson of chapter 303: I can say how the fields get their water.
\end{chapteropening}

\section[\texorpdfstring{\pa{ਨਹਿਰ} (nahir)}{ਨਹਿਰ (nahir)}]{\pa{ਨਹਿਰ} (nahir) — a canal}

\label{lesson:PA-C303-nahir}



\begin{quote}
Before the new one: say the Punjabi for to sprout, then the Punjabi for mustard.
\end{quote}

\subsection*{You'll want to know: \pa{ਨਹਿਰ}}
\textbf{\pa{ਨਹਿਰ}} — \emph{nahir} — \textquotedblleft{}a canal\textquotedblright{}.

\begin{quote}
\pa{ਸਾਡੇ} \pa{ਪਿੰਡ} \pa{ਦੇ} \pa{ਨੇੜੇ} \pa{ਨਹਿਰ} \pa{ਹੈ।}
\end{quote}

\begin{quote}
\emph{sāḍe piṇḍ de neṛe nahir hai.}
\end{quote}

\begin{quote}
There is a canal near our village.
\end{quote}

A feminine noun.

\begin{cousinweb}
\textbf{\pa{ਨਹਿਰ}} (\emph{nahir}) is a loan, usually traced through Persian to Arabic \emph{nahr}, \textquotedblleft{}river\textquotedblright{}; its sense here, \textquotedblleft{}canal\textquotedblright{}, is a shift from that.
\end{cousinweb}

\begin{grammarlens}[title={What you've built}]
One more word for this chapter.
\end{grammarlens}

\subsection*{Your turn}
\begin{itemize}
\raggedright
  \item \emph{Say it:} \emph{nahir}
  \item \emph{Say it:} \emph{nahir}, once more
  \item \emph{Say it:} \emph{sāḍe piṇḍ de neṛe nahir hai.} — There is a canal near our village.
  \item \emph{Recall:} say the Punjabi for to sprout, then the Punjabi for mustard, then say \emph{nahir} again
\end{itemize}

\begin{tcolorbox}[breakable,colback=teal!4,colframe=teal!35!black,title={Before you move on}]
What does \pa{ਨਹਿਰ} mean? (\textquotedblleft{}A canal\textquotedblright{}.) Say it once more.
\end{tcolorbox}
\section[\texorpdfstring{\pa{ਖੂਹ} (khūh)}{ਖੂਹ (khūh)}]{\pa{ਖੂਹ} (khūh) — a well}

\label{lesson:PA-C303-khuh}



\begin{quote}
Before the new one: say the Punjabi for mustard, then the Punjabi for a canal.
\end{quote}

\subsection*{You'll want to know: \pa{ਖੂਹ}}
\textbf{\pa{ਖੂਹ}} — \emph{khūh} — \textquotedblleft{}a well\textquotedblright{}.

\begin{quote}
\pa{ਖੂਹ} \pa{ਦਾ} \pa{ਪਾਣੀ} \pa{ਠੰਢਾ} \pa{ਹੈ।}
\end{quote}

\begin{quote}
\emph{khūh dā pāṇī ṭhaṇḍhā hai.}
\end{quote}

\begin{quote}
The water of the well is cold.
\end{quote}

A masculine noun.

\begin{cousinweb}
\textbf{\pa{ਖੂਹ}} (\emph{khūh}) is inherited from a reconstructed \emph{*khūha}, \textquotedblleft{}pit, well\textquotedblright{}; CDIAL lists Punjabi \emph{khūh}, \textquotedblleft{}well\textquotedblright{}, under it (CDIAL 3912).
\end{cousinweb}

\begin{grammarlens}[title={What you've built}]
One more word for this chapter.
\end{grammarlens}

\subsection*{Your turn}
\begin{itemize}
\raggedright
  \item \emph{Say it:} \emph{khūh}
  \item \emph{Say it:} \emph{khūh}, once more
  \item \emph{Say it:} \emph{khūh dā pāṇī ṭhaṇḍhā hai.} — The water of the well is cold.
  \item \emph{Recall:} say the Punjabi for mustard, then the Punjabi for a canal, then say \emph{khūh} again
  \item \emph{Write it:} copy \textbf{\pa{ਖੂਹ}} once, checking the whole word
\end{itemize}

\begin{tcolorbox}[breakable,colback=teal!4,colframe=teal!35!black,title={Before you move on}]
What does \pa{ਖੂਹ} mean? (\textquotedblleft{}A well\textquotedblright{}.) Say it once more.
\end{tcolorbox}
\section[\texorpdfstring{\pa{ਸਿੰਜਣਾ} (sinjṇā)}{ਸਿੰਜਣਾ (sinjṇā)}]{\pa{ਸਿੰਜਣਾ} (sinjṇā) — to irrigate; to water a field}

\label{lesson:PA-C303-sinjna}



\begin{quote}
Before the new one: say the Punjabi for a canal, then the Punjabi for a well.
\end{quote}

\subsection*{You'll want to know: \pa{ਸਿੰਜਣਾ}}
\textbf{\pa{ਸਿੰਜਣਾ}} — \emph{sinjṇā} — \textquotedblleft{}to irrigate; to water a field\textquotedblright{}.

\begin{quote}
\pa{ਅਸੀਂ} \pa{ਨਹਿਰ} \pa{ਦੇ} \pa{ਪਾਣੀ} \pa{ਨਾਲ} \pa{ਖੇਤ} \pa{ਸਿੰਜਦੇ} \pa{ਹਾਂ।}
\end{quote}

\begin{quote}
\emph{asīṁ nahir de pāṇī nāl khet sinjde hāṁ.}
\end{quote}

\begin{quote}
We irrigate the fields with canal water.
\end{quote}

Past (I …): \pa{ਮੈਂ} \pa{ਖੇਤ} \pa{ਸਿੰਜਿਆ}, \emph{maiṁ khet sinjiā}.

\begin{cousinweb}
\textbf{\pa{ਸਿੰਜਣਾ}} (\emph{sinjṇā}) is inherited from Sanskrit \emph{siñcati}, \textquotedblleft{}pours out, sprinkles\textquotedblright{}; CDIAL lists Punjabi \emph{siñjṇā}, \textquotedblleft{}to water, irrigate\textquotedblright{}, under it (CDIAL 13394).
\end{cousinweb}

\begin{grammarlens}[title={What you've built}]
One more word for this chapter.
\end{grammarlens}

\subsection*{Your turn}
\begin{itemize}
\raggedright
  \item \emph{Say it:} \emph{sinjṇā}
  \item \emph{Say it:} \emph{sinjṇā}, once more
  \item \emph{Say it:} \emph{asīṁ nahir de pāṇī nāl khet sinjde hāṁ.} — We irrigate the fields with canal water.
  \item \emph{Recall:} say the Punjabi for a canal, then the Punjabi for a well, then say \emph{sinjṇā} again
  \item \emph{Write it:} copy \textbf{\pa{ਸਿੰਜਣਾ}} once, checking the whole word
\end{itemize}

\begin{tcolorbox}[breakable,colback=teal!4,colframe=teal!35!black,title={Before you move on}]
What does \pa{ਸਿੰਜਣਾ} mean? (\textquotedblleft{}To irrigate; to water a field\textquotedblright{}.) Say it once more.
\end{tcolorbox}
\section[\texorpdfstring{\pa{ਕਿਸਾਨ} (kisān)}{ਕਿਸਾਨ (kisān)}]{\pa{ਕਿਸਾਨ} (kisān) — a farmer}

\label{lesson:PA-C303-kisan}



\begin{quote}
Before the new one: say the Punjabi for a well, then the Punjabi for to irrigate.
\end{quote}

\subsection*{You'll want to know: \pa{ਕਿਸਾਨ}}
\textbf{\pa{ਕਿਸਾਨ}} — \emph{kisān} — \textquotedblleft{}a farmer\textquotedblright{}.

\begin{quote}
\pa{ਕਿਸਾਨ} \pa{ਸਵੇਰੇ} \pa{ਖੇਤ} \pa{ਜਾਂਦਾ} \pa{ਹੈ।}
\end{quote}

\begin{quote}
\emph{kisān savere khet jāndā hai.}
\end{quote}

\begin{quote}
The farmer goes to the field in the morning.
\end{quote}

A masculine noun.

\begin{cousinweb}
\textbf{\pa{ਕਿਸਾਨ}} (\emph{kisān}) goes back to Sanskrit \emph{kṛṣāṇa}, \textquotedblleft{}ploughing\textquotedblright{}; CDIAL lists Punjabi \emph{kisāṇ}, \textquotedblleft{}farmer\textquotedblright{}, under it and marks it as borrowed rather than inherited, as its \emph{-s-} shows (CDIAL 3447).
\end{cousinweb}

\begin{grammarlens}[title={What you've built}]
One more word for this chapter.
\end{grammarlens}

\subsection*{Your turn}
\begin{itemize}
\raggedright
  \item \emph{Say it:} \emph{kisān}
  \item \emph{Say it:} \emph{kisān}, once more
  \item \emph{Say it:} \emph{kisān savere khet jāndā hai.} — The farmer goes to the field in the morning.
  \item \emph{Recall:} say the Punjabi for a well, then the Punjabi for to irrigate, then say \emph{kisān} again
  \item \emph{Write it:} copy \textbf{\pa{ਕਿਸਾਨ}} once, checking the whole word
\end{itemize}

\begin{tcolorbox}[breakable,colback=teal!4,colframe=teal!35!black,title={Before you move on}]
What does \pa{ਕਿਸਾਨ} mean? (\textquotedblleft{}A farmer\textquotedblright{}.) Say it once more.
\end{tcolorbox}
\section[\texorpdfstring{\pa{ਹਲਟ} (halaṭ)}{ਹਲਟ (halaṭ)}]{\pa{ਹਲਟ} (halaṭ) — a Persian wheel}

\label{lesson:PA-C303-halat}



\begin{quote}
Before the new one: say the Punjabi for to irrigate, then the Punjabi for a farmer.
\end{quote}

\subsection*{You'll want to know: \pa{ਹਲਟ}}
\textbf{\pa{ਹਲਟ}} — \emph{halaṭ} — \textquotedblleft{}a Persian wheel\textquotedblright{}.

\begin{quote}
\pa{ਹਲਟ} \pa{ਖੂਹ} \pa{ਤੋਂ} \pa{ਪਾਣੀ} \pa{ਕੱਢਦਾ} \pa{ਹੈ।}
\end{quote}

\begin{quote}
\emph{halaṭ khūh ton pāṇī kaḍḍhdā hai.}
\end{quote}

\begin{quote}
The Persian wheel draws water from the well.
\end{quote}

A masculine noun.

\begin{cousinweb}
\textbf{\pa{ਹਲਟ}} (\emph{halaṭ}) is inherited from Sanskrit \emph{araghaṭṭa}, \textquotedblleft{}a wheel for raising water\textquotedblright{}; CDIAL lists Punjabi \emph{halhaṭ}, \textquotedblleft{}Persian wheel\textquotedblright{}, under it (CDIAL 596).
\end{cousinweb}

\begin{grammarlens}[title={What you've built}]
One more word for this chapter.
\end{grammarlens}

\subsection*{Your turn}
\begin{itemize}
\raggedright
  \item \emph{Say it:} \emph{halaṭ}
  \item \emph{Say it:} \emph{halaṭ}, once more
  \item \emph{Say it:} \emph{halaṭ khūh ton pāṇī kaḍḍhdā hai.} — The Persian wheel draws water from the well.
  \item \emph{Recall:} say the Punjabi for to irrigate, then the Punjabi for a farmer, then say \emph{halaṭ} again
  \item \emph{Write it:} copy \textbf{\pa{ਹਲਟ}} once, checking the whole word
\end{itemize}

\begin{tcolorbox}[breakable,colback=teal!4,colframe=teal!35!black,title={Before you move on}]
What does \pa{ਹਲਟ} mean? (\textquotedblleft{}A Persian wheel\textquotedblright{}.) Say it once more.
\end{tcolorbox}
